\section{Requisitos del sistema y actores}
\subsection{Actores}
\begin{description}
\item[Ciudadano ficticio:] cuenta autorizada que accede, elige una opción y recibe un comprobante.
\item[Autoridad electoral honesta:] configura la elección, custodia las claves privadas, cierra la urna y ejecuta el escrutinio local.
\item[Atacante externo:] puede intentar enviar peticiones inválidas, manipular datos en tránsito o aprovechar credenciales robadas; no controla la autoridad ni las claves privadas.
\item[Observador:] consulta huellas de sobres cifrados y el manifiesto de resultados firmado.
\end{description}

\subsection{Requisitos funcionales}
\begin{longtable}{p{1.4cm}p{12.5cm}}
\toprule ID & Requisito\\ \midrule\endhead
RF-01 & Crear un padrón \emph{ficticio} de cuentas sin necesidad de servicios oficiales externos.\\
RF-02 & Comprobar usuario y contraseña, sin conservar esta última en texto claro.\\
RF-03 & Mostrar candidatos predefinidos y cifrar la elección localmente antes del envío.\\
RF-04 & Registrar a lo sumo un voto por cuenta mediante una actualización transaccional.\\
RF-05 & Emitir una huella SHA-256 del sobre cifrado sin incluir la candidatura elegida.\\
RF-06 & Registrar una cadena de huellas para detectar modificaciones.\\
RF-07 & Permitir cierre de urna y escrutinio offline de la autoridad honesta.\\
RF-08 & Generar resultados firmados y un listado de auditoría sin identidades.\\
\bottomrule
\end{longtable}

\subsection{Requisitos no funcionales}
\begin{itemize}
\item \concepto{Confidencialidad:} el servidor de votación ordinario no requiere conocer la clave privada para almacenar sobres.
\item \concepto{Integridad:} AES-GCM detecta alteraciones del contenido cifrado durante el escrutinio; SHA-256 y Ed25519 apoyan la integridad de la publicación.
\item \concepto{Autenticación:} se usan contraseñas verificadas con Argon2id y sesiones HTTP protegidas para el entorno local.
\item \concepto{Disponibilidad:} se aceptan limitaciones propias de SQLite y un único proceso de demostración.
\item \concepto{Mantenibilidad:} diseño por módulos, funciones pequeñas, pruebas y documentación.
\end{itemize}

\subsection{Criterios de aceptación}
\begin{enumerate}
\item Una cuenta válida puede emitir exactamente un voto y recibir un comprobante.
\item La misma cuenta no logra registrar un segundo voto aunque dos solicitudes lleguen simultáneamente.
\item Si se modifica un byte de la boleta cifrada, el descifrado autenticado falla.
\item Si un registro almacenado se altera, la auditoría detecta discrepancia.
\item Los resultados tienen una firma Ed25519 verificable.
\item Un observador no encuentra el nombre de usuario ni la candidatura en el registro público de huellas.
\end{enumerate}
